\chapter{Requirements Analysis}
\label{ch:requirements}

This chapter turns the motivation of Chapter~\ref{ch:introduction} into a precise
specification. It identifies the system's actors, enumerates its functional and
non-functional requirements, expresses the key interactions as use cases (with
UML diagrams and detailed descriptions), captures representative user stories,
and presents the domain model as a class diagram. Because \voname{} is a
multi-service, multi-client system, the requirements are organised by
\emph{subsystem} so that they map cleanly onto the architecture and
implementation chapters that follow.

\section{Actors}
\label{sec:actors}

An \emph{actor} is any external entity that interacts with the system. \voname{}
has both human and external-system actors, summarised in
Table~\ref{tab:actors}.

\begin{table}[htbp]
\centering
\caption{System actors.}
\label{tab:actors}
\small
\begin{tabularx}{\textwidth}{L{3.4cm} L{2.2cm} X}
\toprule
\textbf{Actor} & \textbf{Type} & \textbf{Description} \\
\midrule
Guest & Human & An unauthenticated visitor who can register, log in, or accept an invitation. \\
Member & Human & An authenticated user with a \emph{desk} in a workspace; the primary actor. Can chat, join rooms, move around the 2D office, and manage their own profile. \\
Workspace Admin & Human & A member with the \code{ADMIN} role in a workspace; can manage teams, invite and remove members, and edit the office map layout. \\
Workspace Owner & Human & The creator of a workspace, holding the \code{OWNER} role; has all admin powers plus the ability to archive the workspace. \\
Colyseus Server & System & The 2D game server; synchronises avatar positions and pushes presence snapshots to the backend. \\
Agora Platform & System & External real-time media (SFU) that transports room and proximity audio/video. \\
Email/SMTP & System & External mail transport used to deliver transactional e-mails (sign-up, OTP, password reset). \\
\bottomrule
\end{tabularx}
\end{table}

The human actors form a role hierarchy: an Owner is an Admin is a Member is a
(formerly) Guest. Authorisation throughout the system is expressed against the
ranked workspace roles \code{GUEST < MEMBER < ADMIN < OWNER}.

\section{Functional Requirements}
\label{sec:fr}

Functional requirements are grouped by subsystem. Each is given a stable
identifier (\code{FR-\textit{xx}}) for traceability.

\subsection{Identity and Accounts (User Service)}
\begin{description}[leftmargin=1.6cm,style=nextline]
  \item[FR-01] A guest shall register with first name, last name, e-mail,
  phone number, and password; the system shall reject a duplicate e-mail.
  \item[FR-02] A registered user shall log in with e-mail and password and
  receive an authentication token valid for a bounded period.
  \item[FR-03] The system shall hash all passwords and one-time passwords before
  storage.
  \item[FR-04] A user shall request a one-time password (OTP) for password reset
  and verify it within a time limit.
  \item[FR-05] A user shall change their password by supplying the current one.
  \item[FR-06] A user shall view their profile and upload/retrieve a profile
  picture.
  \item[FR-07] The system shall emit an event on sign-up, OTP request, and
  password change so that a confirmation e-mail is sent.
\end{description}

\subsection{Workspaces, Desks, and Teams (Workspace Service)}
\begin{description}[leftmargin=1.6cm,style=nextline]
  \item[FR-08] A member shall create a workspace, becoming its owner and
  receiving an owner desk in the same transaction.
  \item[FR-09] A workspace shall be invite-only; an admin shall invite users by
  e-mail, and an invitation shall expire after a fixed period.
  \item[FR-10] A guest shall accept or decline an invitation via a tokenised
  link; acceptance shall create or reactivate the invitee's desk.
  \item[FR-11] A member shall view the workspace member directory and edit their
  own desk (profile, avatar, status, links, widgets, position).
  \item[FR-12] An admin shall create, rename, and delete teams and assign members
  to them.
  \item[FR-13] An admin shall edit the 2D map layout (tilesets, layers, zones,
  spawn points, map objects) with optimistic-concurrency protection.
  \item[FR-14] The system shall expose an internal session-configuration API that
  provides the game server with the full room boot payload (workspace, active
  desks, active map objects).
  \item[FR-15] The system shall accept avatar presence updates (online status,
  position) from the game server and resolve a member's workspace role on
  request.
\end{description}

\subsection{Messaging (Chat Service)}
\begin{description}[leftmargin=1.6cm,style=nextline]
  \item[FR-16] A member shall create and join group channels within a workspace
  and leave them.
  \item[FR-17] A member shall open a direct-message conversation with another
  user; the system shall guarantee a single DM per user pair.
  \item[FR-18] A member shall send text messages to a channel, DM, or thread and
  receive them in real time.
  \item[FR-19] A member shall reply to a message, mention other users, edit their
  own messages, and soft-delete messages.
  \item[FR-20] A member shall start a thread from a message and reply within it.
  \item[FR-21] The system shall track per-user read position and report unread
  counts for channels and threads.
  \item[FR-22] The system shall broadcast typing indicators.
  \item[FR-23] The system shall guarantee idempotent message delivery so that a
  retried send is not duplicated.
  \item[FR-24] The system shall create, update, and delete \code{ROOM}- and
  workspace-level channels in response to events from the room and workspace
  services.
\end{description}

\subsection{Meeting Rooms and Proximity Voice (Room Service)}
\begin{description}[leftmargin=1.6cm,style=nextline]
  \item[FR-25] A member shall create a voice/video room bound to a text channel,
  and add or remove room members.
  \item[FR-26] A member shall join a room and receive a media access token and
  the current participant list; joins shall respect a capacity limit.
  \item[FR-27] The system shall track live presence in a room (mute, camera,
  screen-share state) and broadcast join/leave/state changes in real time.
  \item[FR-28] The system shall compute proximity-based voice groups from avatar
  positions: users within a zone or within hearing range shall be assigned to a
  shared media channel.
\end{description}

\subsection{2D Virtual Office (SkyOffice)}
\begin{description}[leftmargin=1.6cm,style=nextline]
  \item[FR-29] A user shall appear as an avatar and move around a tile-mapped
  office with keyboard or on-screen joystick, colliding with walls and objects.
  \item[FR-30] A user shall see other users move in real time with smooth motion.
  \item[FR-31] A user shall sit on chairs, and interact with computers
  (screen-share stations) and whiteboards.
  \item[FR-32] Proximity between two avatars shall automatically start and stop a
  peer-to-peer video call; a designated area shall keep a call alive as a private
  meeting space.
  \item[FR-33] A user shall exchange in-world chat messages, shown both in a chat
  panel and as transient speech bubbles.
\end{description}

\subsection{Notifications (Notification Service)}
\begin{description}[leftmargin=1.6cm,style=nextline]
  \item[FR-34] The system shall consume domain events and deliver transactional
  e-mails from versioned HTML templates.
  \item[FR-35] The system shall persist in-app notifications and push them to
  online users in real time, backfilling for offline users.
  \item[FR-36] A user shall list notifications, see an unread count, mark items
  read (individually or all), and delete them.
  \item[FR-37] Notification delivery shall be idempotent with respect to a unique
  event identifier.
\end{description}

\subsection{Cross-cutting Client Requirements}
\begin{description}[leftmargin=1.6cm,style=nextline]
  \item[FR-38] The web client shall provide the full workspace, channel, DM,
  thread, room-call, and notification experience.
  \item[FR-39] The mobile client shall provide the messaging experience (channels,
  DMs, threads) with push notifications across foreground, background, and
  terminated states.
  \item[FR-40] All real-time clients shall authenticate their WebSocket
  connection using a single-use ticket and reconnect automatically.
\end{description}

\section{Non-functional Requirements}
\label{sec:nfr}

\begin{table}[htbp]
\centering
\caption{Non-functional requirements.}
\label{tab:nfr}
\small
\begin{tabularx}{\textwidth}{L{2.2cm} L{3.0cm} X}
\toprule
\textbf{ID} & \textbf{Quality} & \textbf{Requirement} \\
\midrule
NFR-01 & Performance & Avatar movement and chat delivery shall feel instantaneous; movement shall use client-side prediction and interpolation, and state sync shall transmit only changed fields. \\
NFR-02 & Scalability & Media (audio/video/screen) shall never traverse the application servers; it shall use peer-to-peer or an external SFU so that server load is independent of media volume. \\
NFR-03 & Security & Authentication shall be validated once at the gateway; internal endpoints shall be unreachable from the edge; client-supplied identity headers shall be stripped. \\
NFR-04 & Reliability & Event consumers shall retry with exponential backoff and process events idempotently; a service outage shall degrade gracefully rather than cascade. \\
NFR-05 & Consistency & Cross-service state (e.g. room-to-chat channel binding) shall be eventually consistent via durable events, with clients tolerating brief propagation delay. \\
NFR-06 & Maintainability & Each service shall own its data and schema; services shall be independently buildable and deployable. \\
NFR-07 & Portability & Each service shall ship with a container/compose definition for its dependencies to enable local and containerised deployment. \\
NFR-08 & Usability & Clients shall support light and dark themes and responsive layouts across desktop and mobile form factors. \\
NFR-09 & Observability & Services shall expose health endpoints and structured logs for monitoring. \\
\bottomrule
\end{tabularx}
\end{table}

\section{Use-Case Model}
\label{sec:usecases}

Figure~\ref{fig:usecase-overview} gives the top-level use-case diagram, grouping
use cases by subsystem and relating them to the actors of
Section~\ref{sec:actors}.

\begin{figure}[htbp]
\centering
\begin{tikzpicture}[
    every node/.style={font=\small},
    usecase/.style={draw=voPrimary,fill=voPrimary!8,ellipse,minimum width=3.2cm,
                    minimum height=1.2cm,align=center,inner sep=1pt,text width=2.7cm},
    actor/.style={align=center,font=\small},
    sys/.style={draw=voLine,rounded corners=4pt,inner sep=10pt},
]
% ---- actor drawing macro ----
\newcommand{\stick}[3]{% x, y, label
  \begin{scope}[shift={(#1,#2)}]
     \draw[voDark,thick] (0,0.55) circle (0.16);
     \draw[voDark,thick] (0,0.39)--(0,-0.12);
     \draw[voDark,thick] (-0.22,0.20)--(0.22,0.20);
     \draw[voDark,thick] (0,-0.12)--(-0.18,-0.45);
     \draw[voDark,thick] (0,-0.12)--(0.18,-0.45);
     \node[align=center,font=\scriptsize] at (0,-0.75) {#3};
  \end{scope}}

% columns / rows
\def\lx{4.8}\def\rx{8.9}

% actors (left)
\stick{0}{8.0}{Guest}
\stick{0}{4.8}{Member}
\stick{0}{1.4}{Workspace\\Admin/Owner}

% actors (right, systems)
\node[actor,right] at (12.5,7.2) {\textbf{[Agora]}\\media};
\node[actor,right] at (12.5,4.8) {\textbf{[Colyseus]}\\game server};
\node[actor,right] at (12.5,2.4) {\textbf{[SMTP]}\\e-mail};

% system boundary
\node[sys,fit={(2.7,-0.6) (12.0,8.9)},label={[voPrimary,font=\bfseries]above:Virtual Office}] (bnd) {};

% use cases (rows pitch 1.6)
\node[usecase] (uc1) at (\lx,8.0) {Register / Log in};
\node[usecase] (uc2) at (\rx,8.0) {Accept invitation};
\node[usecase] (uc3) at (\lx,6.4) {Manage profile};
\node[usecase] (uc4) at (\rx,6.4) {Manage workspace \& teams};
\node[usecase] (uc5) at (\lx,4.8) {Send / read messages};
\node[usecase] (uc6) at (\rx,4.8) {Threads \& reactions};
\node[usecase] (uc7) at (\lx,3.2) {Join meeting room};
\node[usecase] (uc8) at (\rx,3.2) {Proximity voice/video};
\node[usecase] (uc9) at (\lx,1.6) {Move avatar / interact};
\node[usecase] (uc10) at (\rx,1.6) {Receive notifications};
\node[usecase] (uc11) at (6.85,0.0) {Edit map layout};

% associations - guest (arc the second link over uc1 to avoid the strike-through)
\draw (0.3,8.0)--(uc1.west);
\draw (0.3,8.2) to[out=28,in=150] (uc2.north west);
% member - left column straight
\draw (0.3,4.8)--(uc3.west); \draw (0.3,4.8)--(uc5.west);
\draw (0.3,4.8)--(uc7.west); \draw (0.3,4.8)--(uc9.west);
% member - right column, routed through the row gaps (not through the left ellipses)
\draw (0.3,4.8) to[out=34,in=158] (uc6.north west);
\draw (0.3,4.8) to[out=-34,in=200] (uc8.south west);
\draw (0.3,4.8) to[out=-48,in=196] (uc10.west);
% admin
\draw (0.3,1.4) to[out=35,in=205] (uc4.west);
\draw (0.3,1.4)--(uc11.west);
% systems (secondary actors on the right; routed up the right corridor)
\draw (uc8.east) to[out=0,in=250] (12.4,7.2);           % Agora    <- Proximity voice/video
\draw (uc8.east) to[out=-10,in=200] (12.4,4.8);         % Colyseus <- Proximity (game world)
\draw (uc10.east) to[out=0,in=180] (12.4,2.4);          % SMTP     <- Receive notifications
\end{tikzpicture}
\caption{Top-level use-case diagram for \voname{}.}
\label{fig:usecase-overview}
\end{figure}

\subsection{Detailed Use Cases}

Two representative use cases are specified in full to illustrate the level of
interaction the system supports.

\begin{table}[htbp]
\centering
\caption{Use case UC-05: Send a message to a channel.}
\label{tab:uc-send}
\small
\begin{tabularx}{\textwidth}{L{3.0cm} X}
\toprule
\textbf{Field} & \textbf{Description} \\
\midrule
Use case & Send a message to a channel \\
Actor & Member \\
Preconditions & The member is authenticated and is a member of the target channel; a WebSocket connection is established. \\
Main flow &
1. The member types a message and presses send.
2. The client generates a client message id and publishes the message over STOMP to \route{/app/chat/send}.
3. The chat service verifies channel membership, persists the message (idempotent on the client id), and broadcasts a \code{NEW\_MESSAGE} event to \route{/topic/channel/\{id\}}.
4. All subscribed clients, including the sender, render the message. \\
Alternative flows &
2a. If the socket send fails (mobile), the client falls back to a REST \code{POST} of the same message.
3a. If the member is not a member of the channel, the service returns a \code{NOT\_A\_MEMBER} error on the member's private error queue. \\
Postconditions & The message is stored and visible to all channel members; unread counts update for other members. \\
\bottomrule
\end{tabularx}
\end{table}

\begin{table}[htbp]
\centering
\caption{Use case UC-08: Proximity voice/video in the 2D office.}
\label{tab:uc-prox}
\small
\begin{tabularx}{\textwidth}{L{3.0cm} X}
\toprule
\textbf{Field} & \textbf{Description} \\
\midrule
Use case & Start a proximity call \\
Actor & Member (two members) \\
Preconditions & Both members are in the same office room, have enabled their webcam, and are marked ready to connect. \\
Main flow &
1. Member A moves their avatar toward member B.
2. When the avatars overlap, the client detects the proximity via physics and, to avoid double-calling, the member with the greater session id initiates a PeerJS call.
3. The other member auto-answers; encrypted peer-to-peer audio/video flows directly between them.
4. Each sees the other's video tile appear. \\
Alternative flows &
2a. If either member is not ready or has no webcam, no call is placed.
4a. If the two avatars are inside the designated private-meeting rectangle, the call persists even after they stop overlapping. \\
Postconditions & A live call exists until the members separate (outside the private area), at which point it is torn down. \\
\bottomrule
\end{tabularx}
\end{table}

\section{User Stories}
\label{sec:userstories}

The requirements can also be read as user stories, which capture intent from the
user's perspective:

\begin{itemize}
  \item \emph{As a guest}, I want to register and be invited into a workspace so
  that I can join my team's virtual office.
  \item \emph{As a member}, I want to see my colleagues moving around the office
  so that I know who is present and available.
  \item \emph{As a member}, I want a conversation to start simply by walking up
  to a colleague so that I can ask a quick question without scheduling a meeting.
  \item \emph{As a member}, I want persistent channels, threads, and direct
  messages so that important discussion is not lost when a call ends.
  \item \emph{As a member}, I want to receive a notification on my phone when I am
  mentioned so that I do not miss time-sensitive messages.
  \item \emph{As an admin}, I want to invite and organise members into teams and
  arrange the office map so that the space reflects how my organisation works.
  \item \emph{As an owner}, I want confidence that only authorised users can join
  my workspace's rooms and read its channels.
\end{itemize}

\section{Domain Model}
\label{sec:domain}

Figure~\ref{fig:domain} shows the core domain entities and their relationships,
abstracted across services. It is a conceptual model: in the implementation these
entities are owned by different services and stored in different databases (see
Chapters~\ref{ch:architecture} and~\ref{ch:backend}), but conceptually they form
a single connected model centred on the \emph{Workspace} and the \emph{User}.

\begin{figure}[htbp]
\centering
\begin{tikzpicture}[
  entity/.style={draw=voPrimary,fill=white,rounded corners=2pt,align=left,
                 font=\footnotesize,inner sep=5pt,text width=3.0cm,
                 drop shadow={opacity=0.15,shadow xshift=1pt,shadow yshift=-1pt}},
  rel/.style={voGray,-{Stealth[]},font=\scriptsize},
  lbl/.style={font=\scriptsize,voDark,fill=white,inner sep=1.5pt},
]
\node[entity] (user) at (0,6) {\textbf{User}\\ \hrule id, name\\ email, password\\ phone, verified};
\node[entity] (ws) at (5,6) {\textbf{Workspace}\\ \hrule id, name, slug\\ ownerId, status\\ map layout};
\node[entity] (invite) at (9.7,6.7) {\textbf{Invitation}\\ \hrule email, token\\ status, expiry};
\node[entity] (team) at (9.7,4.2) {\textbf{Team}\\ \hrule id, wsId\\ name};
\node[entity] (desk) at (5,3.3) {\textbf{Desk}\\ \hrule userId, wsId\\ role, status\\ position (x,y)};
\node[entity] (channel) at (0,3.3) {\textbf{Channel}\\ \hrule id, type\\ wsId, members\\ dmKey};
\node[entity] (notif) at (9.9,1.6) {\textbf{Notification}\\ \hrule userId, type\\ body, read};
\node[entity] (message) at (0,0.5) {\textbf{Message}\\ \hrule id, senderId\\ content, threadId\\ replyToId, mentions};
\node[entity] (thread) at (4.0,0.6) {\textbf{Thread}\\ \hrule id, channelId\\ rootMessageId};
\node[entity] (room) at (7.5,0.6) {\textbf{Room}\\ \hrule id, wsId\\ channelId\\ agoraChannel};

% identity / space
\draw[rel] (desk) -- node[lbl]{seats} (user.south east);
\draw[rel] (desk) -- node[lbl,left]{in} (ws);
\draw[rel] (invite) -- node[lbl,right]{to} (ws.east);
\draw[rel] (team) -- node[lbl,above]{groups} (desk.east);
\draw[rel] (channel.north) to[out=90,in=195] node[lbl,pos=0.65]{belongs to} (ws.west);
% chat
\draw[rel] (message) -- node[lbl,left]{in} (channel);
\draw[rel] (message.east) to[out=18,in=162] node[lbl,above]{in} (thread.west);
\draw[rel] (thread.west) to[out=-162,in=-18] node[lbl,below]{roots} (message.east);
% room  (binds arcs high over the thread; membership curves right of the desk)
\draw[rel] (room.north west) to[out=135,in=-28,looseness=1.15]
      node[lbl,pos=0.5]{binds} (channel.south east);
\draw[rel] (room.north) to[out=80,in=-25] node[lbl,right,pos=0.4]{in} (ws.south east);
% notification (routed around the top-right corner, clear of all boxes)
\draw[rel] (notif.east) to[out=0,in=-30] (11.9,7.8) to[out=150,in=35]
      node[lbl,above,pos=0.85]{for} (user.north east);
\end{tikzpicture}
\caption{Conceptual domain model of \voname{} (cross-service view).}
\label{fig:domain}
\end{figure}

The central relationship is that a \emph{User} takes a \emph{Desk} in a
\emph{Workspace}; the desk carries the user's workspace-scoped role, profile, and
avatar position. Everything else --- channels, rooms, teams, invitations --- is
scoped to the workspace, and messaging (messages and threads) hangs off channels.
This model is realised in the following chapters, where each entity's ownership,
storage technology, and API are described in full.
